\chapter{Builds for the Threadweaver}

Theory is necessary. Practice is essential. But a map of the road is not the same as a travelling companion. The following builds are not prescriptions. They are \emph{examples}—casters I have taught, fought alongside, or (in one case) fled from. Their choices reflect their circumstances. Yours will reflect yours.

I have organised these builds by Tier, from novice to master. Each build includes recommended Attributes, Skills, Talents, and signature spells. I have also added a note on playstyle—because the Weave responds to how you think, not only to what you cast.

\section{Tier I: The Cautious Flame (30–40 XP)}

\subsection{Concept}

You have learned enough to be dangerous—to yourself. You cast cantrips reliably. Two- tag spells make you nervous. Dangerous tags are a rumour you have heard from older students. You carry a focus (a stone, a coin, a dried flower) because it helps you remember your intention.

\subsection{Attributes}

\begin{itemize}
    \item \textbf{Wits 3} (for perception and precision)
    \item \textbf{Spirit 2} (for will and resilience)
    \item \textbf{Body 1} (you are not a warrior)
    \item \textbf{Presence 1} (you are not a diplomat)
\end{itemize}

\subsection{Skills}

\begin{itemize}
    \item \textbf{Arcana 2} (your core)
    \item \textbf{Lore 1} (you have read some theory)
    \item \textbf{Insight 1} (you watch witnesses)
\end{itemize}

\subsection{Talents}

\begin{itemize}
    \item \textbf{Spellcraft} (6 XP) – mandatory
    \item \textbf{Practiced Caster} (2 XP) – cantrips cost almost nothing
    \item \textbf{Threadweaver's Patience} (3 XP) – you learn to wait for the right moment
\end{itemize}

\subsection{Signature Spells}

\begin{itemize}
    \item \textbf{Ember Flick} ([BURNING][STRIKE], DV 3) – your only offensive option
    \item \textbf{Warmth} ([HEAL], DV 2) – for Fatigue and minor injuries
    \item \textbf{Hush} ([SILENCE], DV 2) – for escaping notice
\end{itemize}

\subsection{Playstyle}

You stay at the back. You cast one spell per scene, then assess. You do not cast dangerous tags. You save your Boons for re-rolling ones, because a miss at Tier I is a lesson you can afford. Your allies protect you. You protect them from minor inconveniences.

\subsection{Lysandra's Note}

This was my build at thirty XP. I was not impressive. I was \emph{alive}. The Weave respects survival.

\section{Tier II: The Reliable Hand (41–90 XP)}

\subsection{Concept}

You have survived. You know which tags work together and which produce interesting failures. You can cast two- tag spells without sweating. Dangerous tags are no longer rumours; you have used [LEAP] once, and you still have the scar.

\subsection{Attributes}

\begin{itemize}
    \item \textbf{Wits 3} (unchanged; precision matters more than power)
    \item \textbf{Spirit 3} (you have learned to endure backlash)
    \item \textbf{Body 1}
    \item \textbf{Presence 2} (you have started to notice witnesses)
\end{itemize}

\subsection{Skills}

\begin{itemize}
    \item \textbf{Arcana 3} (your core, now reliable)
    \item \textbf{Lore 2} (you have read the forbidden chapters)
    \item \textbf{Insight 1}
    \item \textbf{Stealth 1} (sometimes not casting is the best spell)
\end{itemize}

\subsection{Talents}

\begin{itemize}
    \item \textbf{Spellcraft} (already taken)
    \item \textbf{Controlled Burn} (4 XP) – convert ones to Boons
    \item \textbf{Weave Resilience} (4 XP) – reduce backlash severity
    \item \textbf{Elemental Focus (Fire)} (4 XP) – you have chosen an element
\end{itemize}

\subsection{Signature Spells}

\begin{itemize}
    \item \textbf{Ember Dart} ([BURNING][STRIKE], DV 3) – now with +1 die from Elemental Focus
    \item \textbf{Rising Mist} ([WATER][AREA], DV 3) – concealment and confusion
    \item \textbf{Threshold Step} ([LEAP], DV 4) – your first dangerous tag
    \item \textbf{Still Air} ([DEFEND], DV 2) – for when the enemy gets close
\end{itemize}

\subsection{Playstyle}

You are a generalist. You can hurt, hide, and run. You do not specialise. You keep one Boon in reserve for a re-roll, because you have learned that the Weave's sense of humour is not always kind. You cast dangerous tags only when the alternative is death.

\subsection{Lysandra's Note}

This was my build at sixty XP. I passed my examination. I also set fire to a desk. The Weave accepted the offering.

\section{Tier III: The Specialist (91–150 XP)}

\subsection{Concept}

You have chosen a path. You do not cast every element; you cast \emph{your} element. Fire, or water, or shadow—you have learned its grammar, its moods, its favourite receipts. Dangerous tags are tools, not fears. You carry a ledger of your failures. It is thicker than your spellbook.

\subsection{Attributes}

\begin{itemize}
    \item \textbf{Wits 4} (precision honed)
    \item \textbf{Spirit 4} (backlash is a conversation, not a wound)
    \item \textbf{Body 2} (you have learned to dodge)
    \item \textbf{Presence 2} (witnesses matter)
\end{itemize}

\subsection{Skills}

\begin{itemize}
    \item \textbf{Arcana 4} (your core, now formidable)
    \item \textbf{Lore 3} (you have written marginalia in forbidden texts)
    \item \textbf{Insight 2}
    \item \textbf{Stealth 2}
    \item \textbf{Performance 1} (for veiling in crowds)
\end{itemize}

\subsection{Talents (Fire Specialist)}

\begin{itemize}
    \item \textbf{Spellcraft}
    \item \textbf{Controlled Burn}
    \item \textbf{Weave Resilience}
    \item \textbf{Elemental Focus (Fire)}
    \item \textbf{Elemental Attunement (Fire)} (5 XP)
    \item \textbf{Elemental Rebuke (Fire)} (6 XP)
    \item \textbf{Signature Spell (Ember Dart)} (5 XP)
\end{itemize}

\subsection{Signature Spells}

\begin{itemize}
    \item \textbf{Ember Dart} ([BURNING][STRIKE], DV 3) – now with +2 dice from Focus and Attunement, plus a reroll on failures
    \item \textbf{Cinder Breath} ([BURNING][AREA], DV 3) – crowd control
    \item \textbf{Storm Spark} ([STORM][STRIKE], DV 3) – for when fire is not enough
    \item \textbf{Fold the Wall} ([FOLD][AREA], DV 5) – dangerous, but you have practised
    \item \textbf{Breath Returned} ([REWIND][TOUCH], DV 5) – you have used this once. You will not speak of it.
\end{itemize}

\subsection{Playstyle}

You are a specialist. You know exactly what your element can do and exactly what it cannot. You do not waste spells on problems that a knife could solve. You reserve dangerous tags for the moments that define campaigns. Your allies trust you. They should.

\subsection{Lysandra's Note}

I have met three fire specialists at Tier III. Two are still alive. One is a crater. The Weave does not discriminate; it merely records.

\section{Tier IV–V: The Walking Receipt (151+ XP)}

\subsection{Concept}

You no longer fear backlash. You have learned to read the Weave's mood before you cast. You can feel when a dangerous tag will succeed and when it will merely be \emph{interesting}. You carry no focus; you \emph{are} the focus. Students whisper your name. The Synod has a file. The Weave has a receipt with your signature.

\subsection{Attributes}

\begin{itemize}
    \item \textbf{Wits 5} (you see the Weave's grammar)
    \item \textbf{Spirit 5} (you have argued with backlash and won)
    \item \textbf{Body 3} (you have learned to survive your own failures)
    \item \textbf{Presence 4} (witnesses remember you)
\end{itemize}

\subsection{Skills}

\begin{itemize}
    \item \textbf{Arcana 5} (your core is instinct)
    \item \textbf{Lore 4} (you have written books others fear to read)
    \item \textbf{Insight 3}
    \item \textbf{Stealth 3}
    \item \textbf{Performance 2}
    \item \textbf{Command 2} (for when witnesses need convincing)
\end{itemize}

\subsection{Talents}

\begin{itemize}
    \item All previous, plus:
    \item \textbf{Catastrophe Insurance} (8 XP) – for the day you push too far
    \item \textbf{Echo Recall} (4 XP) – because sometimes the second draft works
    \item \textbf{Safe Harbour} (4 XP) – because you are not alone
\end{itemize}

\subsection{Signature Spells}

You no longer keep a fixed list. You improvise. You cast [GATE] to escape collapsing buildings. You cast [TRANSFORM] to turn an enemy's armour into silk. You cast [REWIND] when a friend falls, and you pay the Weave's price without flinching.

\subsection{Playstyle}

You are a legend. You do not cast spells; you \emph{resolve} situations. The Weave listens when you speak because you have paid your receipts for decades. You are still capable of catastrophic failure. You simply survive it better than most.

\subsection{Lysandra's Note}

I am not Tier V. I have met three Tier V casters in my life. One is dead. One is missing. One is me, on a good day, when the Weave is feeling generous. The rest is practice.

\section{A Final Comparison}

\noindent
\begin{tabular}{|p{0.2\textwidth}|p{0.25\textwidth}|p{0.25\textwidth}|p{0.25\textwidth}|}
\hline
 & \textbf{Tier I} & \textbf{Tier III} & \textbf{Tier V} \\
\hline
\textbf{Cantrips} & Nervous & Confident & Instinctive \\
\hline
\textbf{Two- Tag Spells} & Risky & Routine & Boring \\
\hline
\textbf{Dangerous Tags} & Never & Sometimes & When appropriate \\
\hline
\textbf{Backlash} & Feared & Managed & Welcomed \\
\hline
\textbf{Boons} & Hoarded & Spent & Invested \\
\hline
\textbf{Witnesses} & Avoided & Noticed & Commanded \\
\hline
\end{tabular}

\noindent
The Weave does not care which tier you occupy. It cares about your intention, your grammar, and your willingness to accept the receipt.

Build your caster. Cast your spells. Fail interestingly. Learn.

Then cast again.

\textit{--- Lysandra}
